\PassOptionsToPackage{dvipsnames}{xcolor}
\documentclass[10pt,a4paper,normalphoto]{altacv}

% Géométrie optimisée pour une mise en page moderne
\geometry{left=1cm,right=9cm,marginparwidth=7.2cm,marginparsep=1.0cm,top=0.4cm,bottom=1.0cm,footskip=1.2\baselineskip}

\usepackage[utf8]{inputenc}
\usepackage[T1]{fontenc}
\usepackage[sfdefault,light]{roboto} % Police Roboto moderne
\usepackage{graphicx}
\usepackage{xcolor}
\usepackage{tikz}
\usetikzlibrary{shadings,shadows.blur}
\usepackage[colorlinks=true, urlcolor=DataGreen]{hyperref}
\usepackage[normalem]{ulem}
\usepackage[export]{adjustbox}
\usepackage{fontawesome}
\usepackage[most]{tcolorbox}

% === PALETTE DATA GREEN MODERNE ===
\definecolor{DataGreen}{HTML}{16A085}        % Couleur principale Data Science
\definecolor{DataGreenLight}{HTML}{48C9B0}   % Vert clair pour accents
\definecolor{DataGreenDark}{HTML}{117A65}    % Vert foncé pour texte important
\definecolor{SoftGray}{HTML}{F8F9FA}         % Gris très clair pour fond
\definecolor{ModernGray}{HTML}{2C3E50}       % Gris moderne pour texte
\definecolor{TechWhite}{HTML}{FFFFFF}        % Blanc pur
\definecolor{ShadowGray}{HTML}{BDC3C7}       % Gris pour ombres

% Configuration des couleurs du template
\colorlet{heading}{DataGreenDark}
\colorlet{accent}{DataGreen}
\colorlet{emphasis}{ModernGray}
\colorlet{body}{ModernGray}

\renewcommand{\itemmarker}{{\color{DataGreen}\small\faCircle}}
\renewcommand{\ratingmarker}{\color{DataGreen}\faCircle}

% === TAGS INTERACTIFS AVEC BORDURES COLORÉES ===
\renewcommand{\cvtag}[1]{%
  \tikz[baseline]\node[
    anchor=base,
    draw=DataGreen,
    fill=DataGreenLight!10,
    rounded corners=3pt,
    inner xsep=0.8ex,
    inner ysep=0.6ex,
    text height=1.2ex,
    text depth=0.25ex,
    font=\small\sffamily,
    text=DataGreenDark,
    drop shadow={opacity=0.2,shadow xshift=0.5pt,shadow yshift=-0.5pt}
  ]{#1};%
}

% === DIVIDERS STYLISÉS ===
\renewcommand{\divider}{
  \begin{tikzpicture}
    \draw[DataGreenLight, line width=1.5pt, rounded corners=0.5pt] 
      (0,0) -- (\linewidth,0);
  \end{tikzpicture}
  \medskip
}

% === BANNIÈRE EN-TÊTE AVEC DÉGRADÉ ===
\newcommand{\modernheader}{
  \begin{tikzpicture}[remember picture,overlay]
    \fill[left color=DataGreen, right color=DataGreenLight, opacity=0.1]
      (current page.north west) rectangle ++(\paperwidth,-3cm);
  \end{tikzpicture}
}

% === HIGHLIGHT BOX POUR CERTIFICATIONS ===
\newcommand{\highlightbox}[2]{
  \begin{tcolorbox}[
    enhanced,
    colback=DataGreenLight!8,
    colframe=DataGreen,
    boxrule=1.5pt,
    arc=4pt,
    left=4pt,
    right=4pt,
    top=4pt,
    bottom=4pt,
    drop shadow={opacity=0.15,shadow xshift=1pt,shadow yshift=-1pt}
  ]
    \textbf{\color{DataGreenDark}#1} \\
    \small #2
  \end{tcolorbox}
}

% === SECTION AVEC ENCADRÉS POUR PROJETS ===
\newcommand{\projectbox}[1]{
  \begin{tcolorbox}[
    enhanced,
    colback=SoftGray,
    colframe=DataGreenLight,
    boxrule=1pt,
    arc=3pt,
    left=6pt,
    right=6pt,
    top=6pt,
    bottom=6pt
  ]
    #1
  \end{tcolorbox}
  \vspace{0.3em}
}

% === AMÉLIORATION DE LA HIÉRARCHIE TYPOGRAPHIQUE ===
\renewcommand{\cvsection}[2][]{%
  \bigskip%
  \ifstrequal{#1}{}{}{\marginpar{\vspace*{\dimexpr1pt-\baselineskip}\raggedright\input{#1}}}%
  \begingroup
    \begin{tikzpicture}
      \fill[left color=DataGreen, right color=DataGreenLight, opacity=0.15]
        (0,0) rectangle (\linewidth,0.8cm);
    \end{tikzpicture}
    \vspace{-0.6cm}
    {\color{TechWhite}\LARGE\bfseries\sffamily\MakeUppercase{#2}}
  \endgroup
  \medskip
}

\renewcommand{\cvsubsection}[1]{%
  \smallskip%
  {\color{DataGreenDark}\large\bfseries\sffamily{#1}\par}
  \vspace{0.2em}
  {\color{DataGreenLight}\rule{0.3\linewidth}{1.5pt}\par}
  \medskip
}

% === ÉVÉNEMENTS AVEC STYLE MODERNE ===
\renewcommand{\cvevent}[4]{%
  \begin{tcolorbox}[
    enhanced,
    colback=TechWhite,
    colframe=DataGreenLight,
    boxrule=0.5pt,
    arc=2pt,
    left=8pt,
    right=8pt,
    top=6pt,
    bottom=6pt,
    drop shadow={opacity=0.1,shadow xshift=0.5pt,shadow yshift=-0.5pt}
  ]
    {\large\color{DataGreenDark}\bfseries\sffamily#1\par}
    \smallskip
    {\color{DataGreen}\bfseries#2}\par
    \smallskip
    {\small\color{ModernGray}
      \makebox[0.5\linewidth][l]{\faCalendar \hspace{0.5em}#3}%
      \ifstrequal{#4}{}{}{\makebox[0.5\linewidth][l]{\faMapMarker\hspace{0.5em}#4}}\par}
  \end{tcolorbox}
  \medskip
}

\addbibresource{sample.bib}

\begin{document}

% Ajout de la bannière moderne
\modernheader

\name{\huge{\sffamily\bfseries ANAS SGHIR}}

\tagline{\large{\color{DataGreen}\sffamily\bfseries Data Scientist | Machine Learning Engineer | Product Owner}}

% Photo avec ombre moderne
\photo{2.5cm}{images/Photo.jpg}

\personalinfo{
  \email{anas@itssghir.com}
  \phone{07 58 93 41 75}
  \location{Disponible à partir de septembre 2025}
  \linkedin{\href{https://linkedin.com/in/anassghir}{\color{DataGreen}\uline{LinkedIn}}}
  \github{\href{https://github.com/itssghir}{\color{DataGreen}\uline{GitHub}}}
  \homepage{\href{https://anas.itssghir.com}{\color{DataGreen}\uline{Portfolio}}}
}

\begin{fullwidth}
\makecvheader
\vspace{-0.8em}
\end{fullwidth}

\begin{fullwidth}
\cvsection{\faUser \hspace{0.5em} À PROPOS}
\sffamily Étudiant en Master MIAGE spécialisé en \textbf{\color{DataGreen}Data Science} avec de solides compétences en \textbf{\color{DataGreen}Machine Learning} et \textbf{\color{DataGreen}Statistical Analysis}. Formé aux technologies \textbf{\color{DataGreen}Python}, \textbf{\color{DataGreen}SQL} et \textbf{\color{DataGreen}Cloud Computing} (AWS, Azure) avec une approche pratique des projets Data Science. Expérience en \textbf{\color{DataGreen}Business Intelligence}, \textbf{\color{DataGreen}Predictive Modeling} et \textbf{\color{DataGreen}Data Visualization} (Tableau, Power BI) acquise lors de stages et projets académiques. Passionné par l'innovation technologique, je cherche à appliquer mes compétences techniques dans un environnement professionnel stimulant.
\end{fullwidth}

\cvsection[sidebar_anas_final]{\faBriefcase \hspace{0.5em} EXPÉRIENCES PROFESSIONNELLES}

\cvevent{\textbf{Data Analyst \& Product Owner}}{%
  \includegraphics[height=1em]{images/logo_banque_postale.png} \hspace{0.5em} La Banque Postale}{2025 - En cours}{Toulouse}
\begin{itemize}
\item Data Analysis et Statistical Analysis sur datasets financiers
\item Développement de requêtes SQL et scripts Python
\item Creation de dashboards avec Tableau \& Power BI
\item Predictive Modeling et Data Mining sur projets d'équipe
\item Support à la Product Strategy et reporting métier
\item KPI Tracking et analyse de Performance Metrics
\item Participation aux processus de Data Quality
\end{itemize}

\divider

\cvevent{\textbf{Projet Académique - Application Empreinte Carbone}}{%
  \includegraphics[height=1em]{images/logo_toulouse3.jpg} \hspace{0.5em} Université Toulouse III}{2023 - 2024}{Toulouse}
\begin{itemize}
\item Développement d'une application de calcul d'empreinte carbone (Python, React Native)
\item Reconnaissance caméra et vocale (APIs Google)
\item Gestion complète du projet (coordination, planification, supervision)
\end{itemize}

\divider

\cvevent{\textbf{Data Analyst}}{%
  \includegraphics[height=1em]{images/logo_shl.png} \hspace{0.5em} SHL}{2022 - 2023}{Marseille}
\begin{itemize}
\item Développement web avec Python, JavaScript et SQL
\item Data Analysis et création de rapports
\item Participation aux projets Agile en équipe
\item Statistical Analysis et suivi de métriques
\item Contribution aux activités de Database Management
\end{itemize}

\divider

\cvevent{\textbf{Freelance Développeur Web}}{%
  \includegraphics[height=1em]{images/logo_aix_marseille.jpg} \hspace{0.5em} Indépendant}{2021 - 2022}{Aix-en-Provence}
\begin{itemize}
\item Conception et développement de sites web multisecteurs
\item Travail en équipe interdisciplinaire
\item Compétences en gestion, communication, autonomie
\end{itemize}

% ============= PAGE 2 =============
\clearpage

\cvsection[page2_sidebar]{\faRocket \hspace{0.5em} PROJETS PERSONNELS}

\projectbox{
\cvevent{\textbf{Modèles de Prédiction et Clustering}}{Kaggle - Projet Personnel}{2024}{}
\begin{itemize}
\item Entraînement de modèles ML pour détecter le comportement utilisateurs
\item Clustering K-means et DBSCAN pour segmentation
\item Algorithmes de prédiction avec TensorFlow et Scikit-learn
\item Analyse comportementale et patterns utilisateurs
\item Optimisation des hyperparamètres et validation croisée
\end{itemize}
}

\divider

\projectbox{
\cvevent{\textbf{Chatbot Personnalisé avec API ChatGPT}}{Projet Personnel}{2024}{}
\begin{itemize}
\item Utilisation d'API ChatGPT pour besoins spécifiques
\item Personnalisation du modèle pour domaine métier
\item Interface conversationnelle intelligente
\item Intégration avec base de connaissances entreprise
\item Fine-tuning pour réponses contextuelles
\end{itemize}
}

\divider

\cvevent{\textbf{Dashboards Analytics Tableau \& Power BI}}{Projets Multiples}{2023-2024}{}
\begin{itemize}
\item Création de dashboards interactifs dans Tableau
\item Visualisations avancées avec Power BI
\item KPIs temps réel et reporting automatisé
\item Intégration multi-sources de données
\item Optimisation des performances et UX
\end{itemize}

\divider

\cvevent{\textbf{Site Web de Simulation Phishing}}{Projet Éducatif}{2023}{}
\begin{itemize}
\item Développement d'un site de simulation de phishing
\item Objectif : formation à la cybersécurité
\item Interface réaliste pour tests de sensibilisation
\item Tableaux de bord de résultats et statistiques
\item Approche éthique et éducative
\end{itemize}

\divider

\cvevent{\textbf{Modélisation d'Entreprise TOGAF \& ArchiMate}}{Projet Académique}{2024}{}
\begin{itemize}
\item Modélisation architecturale avec framework TOGAF
\item Diagrammes ArchiMate pour architecture d'entreprise
\item Cartographie des processus métiers
\item Architecture technique et applicative
\item Documentation des flux et dépendances
\end{itemize}

% ============= PAGE 3 =============
\clearpage

\cvsection[page3_sidebar]{\faCogs \hspace{0.5em} TECHNOLOGIES \& RÉALISATIONS}

\cvsubsection{\faGraduationCap \hspace{0.5em} FORMATION \& CERTIFICATIONS}

\cvevent{\textbf{Master 2 MIAGE - Data Science}}{Université Toulouse III Paul Sabatier}{2023 - 2025}{En cours}
\begin{itemize}
\item Spécialisation Machine Learning et Data Engineering
\item Projets académiques en Python, SQL, et Statistical Analysis
\item Mémoire sur l'optimisation d'algorithmes de classification
\item Projets pratiques en Data Analysis et visualisation
\end{itemize}

\divider

\highlightbox{🏆 Core Designer Certificate Dataiku}{
Plateforme Data Science - \textbf{\color{DataGreen}Mars 2025} \\
\small
• Data Pipeline development et Data Preparation \\
• Machine Learning model training avec AutoML \\
• Business Intelligence dashboard creation \\
• Data Visualization et Reporting automatisé
}

\cvsubsection{\faStar \hspace{0.5em} EXPÉRIENCE PROFESSIONNELLE}

\cvevent{\textbf{Data Analyst \& Product Owner}}{La Banque Postale}{2025}{Stage/Alternance}
\begin{itemize}
\item Data Analysis sur datasets financiers avec Python et SQL
\item Database Management et Query Optimization (amélioration 15\%)
\item Dashboard creation avec Tableau pour Business Intelligence
\item Participation à des projets Agile/Scrum en équipe
\end{itemize}

\divider

\cvevent{\textbf{Data Analyst}}{SHL - Projet d'équipe}{2022 - 2023}{Mission académique}
\begin{itemize}
\item Statistical Analysis et Data Mining sur données clients
\item Web Development avec JavaScript et Database Integration
\item Project Management et Scrum methodology application
\item Data Visualization avec Excel et Power BI
\end{itemize}

\cvsubsection{\faAward \hspace{0.5em} PROJETS \& CONTRIBUTIONS}

\cvevent{\textbf{Projets Data Science}}{Portfolio Personnel}{2023 - 2025}{Développement continu}
\begin{itemize}
\item Machine Learning projects avec Scikit-learn et TensorFlow
\item Data Analysis et Predictive Modeling sur datasets publics
\item GitHub contributions à des projets open source
\item Technical writing et documentation de projets
\end{itemize}

\divider

\cvevent{\textbf{Développement Personnel \& Apprentissage}}{Formation Continue}{2024 - 2025}{Auto-formation}
\begin{itemize}
\item Cours en ligne sur les technologies Cloud (AWS, Azure)
\item Veille technologique sur les nouvelles méthodes ML
\item Participation à des forums et communautés tech
\item Documentation de projets personnels sur GitHub
\end{itemize}

\end{document} 